\documentclass[12pt,a4paper]{article}
\usepackage[a4paper,margin=2.2cm]{geometry}
\usepackage[T1]{fontenc}
\usepackage[utf8]{inputenc}
\usepackage{newtxtext}
\usepackage{setspace}
\usepackage{xcolor}
\usepackage{tcolorbox}
\usepackage{titlesec}
\usepackage{fancyhdr}
\usepackage{parskip}

\definecolor{luhblue}{RGB}{21,101,192}
\definecolor{soft}{RGB}{245,247,250}

\setstretch{1.25}
\pagestyle{fancy}\fancyhf{}
\lhead{\small SFEM oral examination --- speaking script}
\rhead{\small\thepage}
\renewcommand{\headrulewidth}{0.4pt}

\newcounter{slidenum}
\newcommand{\slide}[3]{%
  \bigskip
  \noindent\colorbox{luhblue}{\parbox{\dimexpr\linewidth-2\fboxsep}{%
    \color{white}\bfseries\small SLIDE #1 \quad|\quad #2 \hfill #3}}\par
  \nopagebreak\vspace{0.3em}}

\newcommand{\num}[1]{\textbf{#1}}
\newcommand{\pron}[2]{#1\,{\scriptsize\itshape[#2]}}

\begin{document}

\begin{center}
{\LARGE\bfseries Speaking Script}\\[0.3em]
{\large Reliability Assessment of the JAXA H3 Fairing}\\[0.2em]
{\normalsize Keisuke Nishioka \quad--\quad 28 August 2026, 9:30 \quad--\quad Room 226, Geb.\ 3403}
\end{center}

\vspace{0.5em}
\begin{tcolorbox}[colback=soft,colframe=luhblue,boxrule=0.6pt,arc=2pt]
{\bfseries Pronunciation --- the ones that trip people up}\\[0.3em]
\pron{Hermite}{er-MEET, not her-mite} \quad
\pron{stochastic}{sto-KAS-tik} \quad
\pron{quadrature}{KWOD-ra-cher} \quad
\pron{orthogonal}{or-THOG-o-nal}\\
\pron{Karhunen--Lo\`{e}ve}{kar-HOO-nen lo-EV} \quad
\pron{Smolyak}{SMOL-yak} \quad
\pron{Sobol}{so-BOL} \quad
\pron{eigenvalue}{EYE-gen-value}\\
\pron{surrogate}{SUR-ro-gat} \quad
\pron{cohesive}{ko-HEE-siv} \quad
\pron{aleatory}{AY-lee-a-tory} \quad
\pron{epistemic}{ep-i-STEE-mik} \quad
\pron{anisotropic}{an-eye-so-TROP-ic}

\smallskip
{\bfseries Pace.} About \num{1{,}900} words for \num{15} minutes --- roughly \num{125}
words a minute. That is deliberately unhurried. If you finish at 13 minutes, good:
the ten minutes of questions is where the marks are.
\end{tcolorbox}

%----------------------------------------------------------
\slide{1}{Title}{0:00 -- 0:20}

Good morning. My name is Keisuke Nishioka, and I will present my project for the
Stochastic Finite Element Methods course.

The question is simple to state. The material properties of a composite fairing vary
from manufacturing. How does that variability propagate into the structural response,
and is the structure reliable?

%----------------------------------------------------------
\slide{2}{Motivation}{0:20 -- 1:05}
2
The H3 payload fairing is a carbon-fibre skin bonded to an aluminium honeycomb core.

The properties of both the laminate and the adhesive are inherently variable ---
fibre volume fraction, cure cycle, adhesive film thickness. These uncertainties
propagate into stress, damage and deflection, and therefore into reliability.

A deterministic analysis at nominal properties cannot tell us how much of the computed
margin is real, and how much is an artefact of the particular values we assumed.

My goal is to quantify how five uncertain material parameters propagate to the
structural response --- treating the existing Abaqus model as a black box.

%----------------------------------------------------------
\slide{3}{The structure}{1:05 -- 1:45}

This is the structure. On the left, the H3 launch vehicle, with the fairing at the nose.

The wall is a sandwich: a carbon-fibre skin, an aluminium-honeycomb core, a second skin,
bonded with an adhesive that I model as a cohesive zone. Diameter \num{5.2} metres,
skins \num{0.6} to \num{1.4} millimetres, core \num{8} millimetres.

Variability lives in the skin stiffness and in the bond. Those are my five uncertain
constants.

%----------------------------------------------------------
\slide{4}{Approach at a glance}{1:45 -- 2:25}
4
Here is the whole approach on one slide.

On the left, the five uncertain material parameters, each given as a distribution.
In the middle, the deterministic finite-element model, which I evaluate at selected
points. On the right, the two things I recover: the probability distribution of the
peak stress, and the Sobol sensitivity indices.

Everything that follows is detail on these three stages.

%----------------------------------------------------------
\slide{5}{Method --- Polynomial Chaos Expansion}{2:25 -- 3:25}

For a response $Y$ depending on a random input vector, polynomial chaos approximates $Y$
as a finite sum of coefficients times orthogonal polynomial basis functions. For Gaussian
inputs those are Hermite polynomials --- the basis is fixed by the input measure, not
chosen arbitrarily.

The crucial point is not the polynomial form. It is the \textbf{orthogonality}.

Because the basis is orthogonal, the coefficients are the spectral projections of the
response. The mean is the zeroth coefficient. The variance is the sum of the remaining
squared coefficients. And the Sobol indices follow analytically, as partial sums of those
squared coefficients --- with \textbf{no additional model evaluations}.

That is the decisive advantage over a generic regression surrogate.

%----------------------------------------------------------
\slide{6}{Method --- Smolyak sparse quadrature}{3:25 -- 4:15}

To build the surrogate I evaluate the model at a set of design points.

A full tensor grid needs $n$ to the power $d$ points: \num{243} runs at degree two,
\num{1024} at degree three. Each run produces about \num{293} megabytes, so the
sample count is the binding constraint.

The Smolyak sparse grid retains exactness on the total-degree space at a fraction of that:
\num{66} points at degree two, \num{286} at degree three.

One practical note. Sparse rules generally produce negative weights, so I do not project
--- I fit the coefficients by \textbf{least-squares regression}. With \num{66} points
against \num{21} unknowns the system is over-determined, which also permits
leave-one-out validation.

%----------------------------------------------------------
\slide{7}{Five uncertain parameters}{4:15 -- 5:00}

These are the five parameters.

Two are in the skin: the fibre-direction modulus $E_1$, at \num{146} gigapascals with a
coefficient of variation of \num{5} percent, and the shear modulus $G_{12}$ at
\num{10} percent.

Three are in the cohesive law: normal stiffness, mode-one fracture energy, and strength,
at \num{15} to \num{20} percent. That ordering is the usual one --- the
fibre-dominated modulus is the most tightly controlled.

I model them as truncated normals between \num{50} and \num{150} percent of nominal,
because stiffnesses and fracture energies are strictly positive.

%----------------------------------------------------------
\slide{8}{FE model and quantities of interest}{5:00 -- 5:45}
8
The deterministic model is an Abaqus/Standard model of a representative panel.

The skin is an orthotropic elastic lamina. The honeycomb core is a homogenised
linear-elastic continuum. The adhesive interface is a cohesive zone model with a
Benzeggagh--Kenane mixed-mode damage law.

The load is applied in two static steps: a pressure, then an incremental load.

There are three quantities of interest: the peak von Mises stress in the skin, with a
limit of \num{1200} megapascals; the maximum cohesive damage variable, SDEG; and
the maximum displacement.

The pipeline is fully non-intrusive. For each design point I patch the material cards in
the input file, run Abaqus, and read the results out of the output database.

%----------------------------------------------------------
\slide{9}{One black-box evaluation}{5:45 -- 6:20}

To make concrete what one evaluation looks like --- this is a von Mises stress field
from an Abaqus output database of the full fairing.

Each quadrature node corresponds to exactly one such deterministic solve. The stress
concentrates around the access-door cutout.

This global model is shown for context; the stochastic study itself wraps a
representative sandwich panel of this structure.

%----------------------------------------------------------
\slide{10}{Result 1 --- convergence}{6:20 -- 7:05}
10
Now the results.

The response statistics are essentially identical between degree two and degree three.
The peak stress is \num{209.7} megapascals in both cases, with a standard deviation
of \num{2.375} against \num{2.378}. They agree to four significant figures.

That is more than a convergence check. It quantifies the curvature of the response
surface: if the response depended appreciably on higher-order terms, the additional
\num{220} simulations would have moved the standard deviation. They did not.

So degree two is not an economy measure. It is the \textbf{correct model order} for this
response.

And notably, the cohesive damage variable is identically zero in all
\num{352} simulations.

%----------------------------------------------------------
\slide{11}{Result 2 --- Sobol indices}{7:05 -- 7:55}

The sensitivity analysis gives a strikingly anisotropic picture.

The fibre-direction modulus alone explains \num{99.2} percent of the variance in
peak stress, and \num{99.95} percent of the variance in displacement.

No other parameter reaches even the one-percent level. The shear modulus is the only
secondary contributor, and it accounts for under one percent of the stress variance ---
despite carrying twice the coefficient of variation of $E_1$.

The three cohesive parameters register indices numerically indistinguishable from zero.

And the first-order and total-order indices are essentially equal. That tells us the
inputs act additively, with no resolvable interactions --- the signature of a response
that is essentially linear over the sampled range.

%----------------------------------------------------------
\slide{12}{Why does $E_1$ dominate?}{7:55 -- 8:50}

Why does one parameter dominate so completely?

To first order, the output variance is the sum of squared sensitivity times input
variance. Each parameter contributes through the \textbf{product} of its sensitivity and
its scatter --- not either alone.

The panel carries the pressure by membrane action, and for membrane action both the peak
stress and the deflection are set by the axial laminate stiffness. So the response is
close to linear in $E_1$, which makes the derivative large --- and $E_1$ still carries
five percent scatter.

The shear modulus is the counter-example: its scatter is twice as large, but the
membrane-dominated kinematics make its derivative small.

The ranking is governed not by which parameter is most uncertain, but by which uncertainty
the structure amplifies.

%----------------------------------------------------------
\slide{13}{Result 3 --- PCE versus a neural network}{8:50 -- 9:45}

How does polynomial chaos compare with a neural network on the same data?

At \num{66} points both are accurate to about ten to the minus three megapascals ---
and on this leave-one-out metric the network is in fact marginally lower.

So accuracy at the design point is not the distinction. \textbf{Stability} is.

Enlarging the training set from \num{66} to \num{286} points makes the network's
error \emph{grow} by a factor of twenty-five: a network with several thousand weights
overfits when the map is this smooth. Degree two holds; degree three is slightly worse,
which is mild over-parameterisation.

And the expansion returns the Sobol indices analytically. The network would need a
separate Monte Carlo study.

%----------------------------------------------------------
\slide{14}{Reliability}{9:45 -- 10:40}

Monte Carlo of a million realisations gives a failure probability of zero and an infinite
reliability index. The mean stress sits about \num{416} standard deviations below the
threshold.

I want to state this carefully. A failure probability of \emph{exactly} zero should be read
as \textbf{below the resolution of the analysis}, not as a guarantee. The Monte Carlo
resolves only to about ten to the minus six, and the surrogate is valid only over the
truncated input domain.

There is also a consequence of the damage being identically zero. The three cohesive
parameters are \textbf{structurally unidentifiable}: they never move any output, so no
sampling could give them a non-zero index. The five-dimensional problem collapses, at this
load, to effectively one dimension.

%----------------------------------------------------------
\slide{15}{Extension T9 --- non-Gaussian random field}{10:40 -- 11:50}

The first extension relaxes the assumption that $E_1$ is a single number.

I represent it as a random field, discretised by Karhunen--Lo\`{e}ve expansion with an
anisotropic exponential kernel.

A Gaussian field admits negative stiffness, which is not admissible. So I map it through a
memoryless transform to obtain a \textbf{translation field} with a lognormal marginal ---
positive by construction.

The transform preserves rank correlation but \emph{not} linear correlation. At the five
percent scatter of $E_1$ that distortion is only three times ten to the minus four, so
here the practical gain is positivity rather than changed statistics. At sixty percent it
reaches nearly four percent --- so non-Gaussian treatment matters for the interface
parameters, not for $E_1$.

Separately, spatial averaging cuts the effective panel modulus scatter from
\num{7300} to about \num{4700} megapascals.

I want to be careful here. That is an \textbf{input-side} result. For responses governed
by average stiffness it means the scalar model is conservative. But the peak stress is a
local extreme, and a heterogeneous field creates gradients that can raise the peak even as
the average varies less. Settling that needs mapped realisations run through the solver,
which I have not done.

%----------------------------------------------------------
\slide{16}{Extension T13 --- surrogate comparison}{11:50 -- 12:55}

The second extension compares polynomial chaos, a Gaussian process and a neural network
over training sets from ten to \num{286} points.

I want to be upfront about what this can and cannot show. The reference response is itself
a degree-two polynomial, so a degree-two expansion contains it exactly. The
ten-to-the-minus-thirteen figure measures \textbf{exact representability} --- not
generalisation error.

The factor that \emph{is} meaningful is the roughly five hundred between the two
data-driven surrogates.

What the comparison measures is the value of a correct prior under a severe data budget.
Had the reference contained a damage bifurcation, the ranking could reverse.

One more point: the Gaussian process supplies an epistemic error estimate the expansion
does not --- \num{0.008} megapascals against an aleatory scatter of \num{2.37}, that
is \num{0.3} percent.

%----------------------------------------------------------
\slide{17}{Extension T6 --- Bayesian identification}{12:55 -- 13:55}

The third extension turns the problem around.

$E_1$ explains ninety-nine percent of the variance, so among the five parameters it is
also the most identifiable from response measurements.

I set this up as a conjugate Gaussian inverse problem, using the linear part of the
expansion as the likelihood. Prior Gaussian, likelihood linear-Gaussian, so the posterior
is closed-form --- no sampling and no new Abaqus runs. Each point uses the mean of $k$
independent measurements, so the estimator is consistent.

The posterior standard deviation falls from \num{7300} megapascals to \num{2582}
at five measurements, \num{1885} at ten --- a \num{74} percent reduction --- and
\num{1356} at twenty.

I should be precise about what this reduces. It is \textbf{epistemic} uncertainty. The
aleatory scatter between batches is untouched, and since the failure probability is
already zero at a four-hundred-sixteen sigma margin, this does \emph{not} improve the
reliability assessment. Its value is on the manufacturing side --- the dominant parameter
is also the one cheap measurements can pin down, which is what tolerance allocation
needs.

%----------------------------------------------------------
\slide{18}{Conclusions}{13:55 -- 15:00}

To conclude.

A degree-two expansion from \num{66} simulations captures the full response statistics,
and I argue this is the correct model order --- not an economy measure.

The actionable message is that the fibre-direction modulus controls more than ninety-nine
percent of the structural variability. Tightening its quality control is by far the most
effective lever.

The extensions sharpen that picture: the translation field secures positivity, the
surrogate comparison shows a correct prior beats flexibility at small sample sizes, and
the Bayesian update shows the dominant parameter is also the identifiable one.

I will end on the limitations. The zero failure probability is resolution-limited. The
inputs are assumed independent. Only a single panel is modelled. And the random-field
realisations were never propagated through the solver.

The natural continuation is a load level that activates the cohesive zone, under a
correlated input model.

Thank you very much for your attention. I am happy to take questions.

%----------------------------------------------------------
\vspace{1em}
\begin{tcolorbox}[colback=soft,colframe=luhblue,boxrule=0.6pt,arc=2pt]
{\bfseries When you do not know the answer}\\[0.4em]
``I don't know. I would have to check that --- my first guess would be \ldots, but I'd want
to verify it before claiming anything.''\\[0.3em]
``That's outside what I investigated. If I were to look into it, I would start by \ldots''\\[0.3em]
``I can answer part of that. \ldots\ The rest I'm not sure about.''\\[0.3em]
``Yes --- that's a limitation I'm aware of. The proper treatment would be \ldots''\\[0.3em]
``You're right --- I hadn't considered that.''\\[0.5em]
{\bfseries Buy time:} ``That's a good question --- let me think for a moment.''\\
{\bfseries Clarify:} ``Do you mean \ldots\ or \ldots?''\\[0.5em]
{\bfseries Never} answer confidently when unsure. A confident wrong answer costs more than
an honest gap.
\end{tcolorbox}

\end{document}
